\subsection{Información y Avisos \hfill Nixon Piguave}
\label{subsec:avisos-nixon}

\subsubsection{Alcance y responsabilidad}
El servicio produce información derivada de CLAB para consulta: reportes
de uso (RF-13), avisos internos (RF-15) y vista central de auditoría
(RF-14). Agrupa módulos que comparten proyecciones y recuperación de
resultados; no cambia la ocupación de Agenda, los permisos de Identidad ni
el estado técnico de Laboratorios. Su motivo de cambio son las reglas para
presentar, distribuir y actualizar esos resultados. Si sus módulos
necesitaran ciclos de despliegue o escalado distintos, se revisaría este
límite en otro ADR.

\subsubsection{Datos de propiedad exclusiva}
La base PostgreSQL \texttt{clab\_avisos} contiene proyecciones de reservas,
asistencia, laboratorios, incidencias y catálogos académicos; notificaciones,
destinatarios y estados de lectura; copias centrales de auditoría; auditoría
local, inbox, estado de sincronización y outbox. Cada servicio conserva
su base y credencial propias. No existen claves foráneas ni consultas SQL
entre servicios. Los identificadores externos son referencias lógicas.

Las proyecciones no sustituyen la verdad transaccional. El reporte conserva
la instantánea de carrera, asignatura y periodo asociada a la reserva para
que un cambio posterior del catálogo no altere su significado histórico.
No se copian contraseñas ni tokens. El estado de permisos proyectado no se
utiliza para autorizar solicitudes.

\subsubsection{Componentes internos}
La vista C3 descompone un contenedor por página, siguiendo el modelo
C4\cite{brownC4Componentes}. La página de Información y Avisos contiene:
\begin{description}
\item[API REST y control de acceso.] Valida filtros, identidad y permiso
vigente; dirige la solicitud al módulo adecuado y restringe los avisos al
destinatario autenticado.
\item[Consultas de uso y exportación.] Consulta proyecciones locales por
carrera y periodo; entrega vista paginada o CSV y la marca de actualización.
Distingue horas reservadas de asistencia realmente registrada.
\item[Gestión de avisos.] Lista avisos por destinatario y actualiza su
lectura idempotente. Guarda auditoría local y outbox en la misma transacción.
\item[Consulta de auditoría.] Recupera las copias centrales por usuario,
origen e intervalo. No modifica las acciones de los servicios fuente.
\item[Consumidores y proyector.] Validan el esquema del evento,
deduplican por consumidor e identificador y actualizan proyección, aviso o
auditoría en una transacción local. Dan ACK después del commit.
\item[Repositorios.] Acceden exclusivamente a la base propia, incluidos
inbox, versiones y marcas de sincronización.
\item[Publicador de outbox.] Publica las acciones locales pendientes con
confirmación; no vuelve a emitir eventos consumidos, evitando bucles.
\end{description}

\subsubsection{API y canales sincrónicos}
Los clientes web y móvil acceden por el Gateway. Se proponen estas operaciones:
\begin{itemize}
\item \texttt{GET /api/reportes/uso}.
\item \texttt{GET /api/reportes/uso/exportacion}.
\item \texttt{GET /api/notificaciones}.
\item \texttt{PATCH /api/notificaciones/\{id\}/lectura}.
\item \texttt{GET /api/auditoria}.
\end{itemize}
La API usa REST/JSON sobre HTTPS porque
el usuario necesita el resultado inmediato y el PFC ya utiliza HTTP/JSON.
No se dispone de mediciones auténticas de PE-U1 para justificar gRPC.

Gateway valida el token y enruta; el servicio comprueba autorización vigente
en Identidad por REST y aplica el control del recurso. Un JWT todavía
vigente no acredita por sí solo que la sesión o el permiso no fueron
revocados. Se conserva el fallo cerrado de RNF-03 y ADR-001: Identidad
inaccesible produce HTTP 503, sin ejecutar operaciones protegidas. El
timeout total de 800 ms para esta comprobación desde Avisos es un presupuesto
propuesto, no un resultado medido ni un presupuesto para Agenda.

\subsubsection{Eventos y mecanismo asincrónico}
Se propone RabbitMQ, exchange topic durable \texttt{clab.eventos.v1},
AMQP/TLS y tres colas durables independientes para proyecciones,
notificaciones y auditoría. Reservas publica confirmaciones, cancelaciones,
conflictos, asistencia y cambios de bloqueo efectivo; Laboratorios,
catálogo, equipos e incidencias; Identidad, cambios académicos y avisos de
cuenta. Todos publican sus acciones auditadas. Avisos produce únicamente
las acciones de sus propias mutaciones. Los payloads y productores deben
validarse entre responsables.

Cada origen guarda negocio, auditoría local y outbox en una transacción.
La publicación posterior emplea mensaje persistente, publisher confirms y
la opción \texttt{mandatory}; un mensaje devuelto sin ruta no se marca
entregado. Los confirms acreditan recepción por el broker, no finalización
del consumidor\cite{nixonRabbitConfirms}. El flujo de negocio no espera a
que Avisos procese el cambio. RNF-06 exige continuidad con Avisos caído y
RNF-07 conserva la auditoría fuente\cite{nixonPfcRNF}. RNF-05 permite
instalar los componentes en la institución sin nube pública obligatoria.

ADR-003 registra esta evolución porque ADR-002 proponía consulta periódica
de cambios mediante REST\cite{nixonPfcADR002}. Kafka sigue siendo una
alternativa de registro retenido; los requisitos revisados no aportan carga
ni reproducción de historial que obliguen a elegirlo. No se afirma una
ventaja experimental de RabbitMQ ni cumplimiento medido de RNF-01/02.

\subsubsection{Fallos, consistencia y recuperación}
La entrega es al menos una vez. Inbox y efectos se guardan juntos; si se
pierde el ACK, la reentrega se deduplica. Las proyecciones usan instantáneas
completas y versión por agregado, de modo que una versión vieja no revierte
una nueva. Auditoría y avisos conservan cada acción única aunque llegue
fuera de orden. No existe orden global entre servicios.

Los errores transitorios tienen reintentos acotados; al quinto fallo se
conserva una copia en la cola de errores antes de dar ACK al original.
Si falla esa copia, no se confirma el original. No se utiliza reencolado
infinito. El detalle de recuperación complementa los mecanismos de errores
del broker\cite{nixonRabbitDLX}.

Si Avisos o RabbitMQ fallan, los orígenes acumulan cambios mientras exista
capacidad; se supervisan disco y backlog. Los reportes muestran la última
marca de agua por origen verificada con un manifiesto de eventos relevantes
por consumidor; no confunden un salto de secuencia por filtrado con pérdida.
Sin reconciliar, indican sincronización no verificada. Si la base local
no responde, se devuelve 503; una proyección no inicializada no aparenta un
reporte vacío. La reconstrucción exige instantáneas consistentes con
secuencia de corte y cambios retenidos; auditoría necesita el archivo de
acciones completo. Retención inicial propuesta: al menos siete días,
supeditada a respaldo y cursores confirmados, con reconciliación si se supera.

La mensajería no establece el bloqueo de agenda: Reservas conserva su
autoridad y el protocolo I13 coordina la barrera operativa antes de anunciar
el cierre del catálogo. Estos contratos, la retención y los detalles de
autorización se presentan para validar con el equipo. Los escenarios de
caída, duplicación, revocación, orden y restauración están documentados en
\path{docs/contratos/eventos-nixon.md}; todavía no se ejecutaron.
